\section{What the model tests}
\label{sec:discussion}

The finite example separates three kinds of restriction that a physical
proposal must not confuse. Admissibility excludes structures that are not
in $\Omega$. Interference can give zero probability to an alternative
assembled from allowed histories. Decoherence determines when those
alternatives can be treated with additive probabilities. None of these
restrictions substitutes for the other two.

\subsection{The weighting must constrain predictions}

The same possibility space supports every row of
Table~\ref{tab:toyoutputs}. Its unlabelled topology is unchanged as well.
Thus admissibility and topology do not select the predictions in this
example. This is not an argument that structure could never constrain a
weighting. It identifies information that the particular relational
structure does not contain.

Allowing an arbitrary positive kernel for each experimental setting would
leave the proposal without a predictive weighting law. Even diagonal
kernels allow arbitrary probability assignments to the four histories.
The example obtains a restriction by holding the mixing operations fixed
and prescribing how phase and record overlap enter their composition.
Explaining that prescription structurally would be a further result.
Renaming its matrices as local rules would not be such an explanation.

The finite example also distinguishes quantum interference from an
ontological assertion that one history is randomly selected. It calculates
probabilities for outputs, but supplies no mechanism selecting one outcome
and does not decide between interpretations of quantum mechanics.

\subsection{Failure criteria}

For the supplied two-path prescription, the phase dependence must obey
Equation~\ref{eq:toyfringes}, visibility must agree with the record overlap,
and the diagonal approximation must have the discrepancy in
Equation~\ref{eq:classicalerror}. A comparison requires independently
specified preparation, mixing operations, record coupling, and detector
conditions. Adjusting these separately for every observed probability
would not test the prescription. Losses would require additional
alternatives or an explicitly conditioned experiment.

These are tests of a standard quantum toy model, not distinctive predictions
of a new spacetime theory. Nor does a two-output interference curve by
itself rule out all classical models with setting-dependent probabilities.
The broader target still requires compatible predictions across different
measurement operations, including genuinely quantum correlations.

A spacetime candidate would face additional failure criteria. Its record
predictions must survive the proposed removal of the regulator. The same
microscopic prescription must account for both geometric and matter
observations without independent fits. Any claimed robustness must persist
over stated changes to microscopic structure, rather than only over the
two continuous parameters used here. Failure on any of these conditions
would reject that candidate, not be repaired by the generality of the
framework.

\subsection{What remains unexplained}

The example shows how a weighted history description can retain
interference while allowing particular coarse records to have ordinary
probabilities. It also shows where the quantum content was supplied.
There is no derived metric, physical clock, gravitational dynamics,
quantum field theory, or stable large-scale geometry in the calculation.
The mixing order is imposed, and the auxiliary record space lies outside
the bare relational description.

The remaining structural problem is therefore specific. A candidate must
explain a restricted composition and weighting prescription, give an
internal account of records, and identify a regime with a common geometric
and quantum description. The calculations here provide checks that such
a candidate should pass. They do not establish that a suitable candidate
exists.
